% GENERATED FILE. Edit canonical lessons, then run npm run generate:books.
% canonical-source-hash: fnv1a64:75c2ce236b5327db
% canonical-lessons: MR-C315-pratyaksadarsi, MR-C315-mahitipat, MR-C315-ank, MR-C315-nondavne, MR-C315-bhatakne

\chapter{Eyewitness Accounts}
\label{ch:mr-eyewitness-accounts}
\hlchaptermodality{315}

\begin{chapteropening}
\textbf{By the end of this chapter:} \emph{I can say who saw something and how it was put on record.}

Complete the last lesson of chapter 315: I can say who saw something and how it was put on record.
\end{chapteropening}

\section[\texorpdfstring{\mr{प्रत्यक्षदर्शी} (pratyakṣadarśī)}{प्रत्यक्षदर्शी (pratyakṣadarśī)}]{\mr{प्रत्यक्षदर्शी} (pratyakṣadarśī) — an eyewitness}

\label{lesson:MR-C315-pratyaksadarsi}



\begin{quote}
Before the new one: say the Marathi for a title, then the Marathi for a front page.
\end{quote}

\subsection*{You'll want to know: \mr{प्रत्यक्षदर्शी}}
\textbf{\mr{प्रत्यक्षदर्शी}} — \emph{pratyakṣadarśī} — \textquotedblleft{}an eyewitness\textquotedblright{}.

\begin{quote}
\mr{एका} \mr{प्रत्यक्षदर्शीने} \mr{पोलिसांना} \mr{सगळं} \mr{सांगितलं}.
\end{quote}

\begin{quote}
\emph{ekā pratyakṣadarśīne polisānnā sagḷaṁ sāṅgitlaṁ.}
\end{quote}

\begin{quote}
An eyewitness told the police everything.
\end{quote}

Masculine.

\mr{प्रत्यक्ष}, before the eyes, + -\mr{दर्शी}, one who sees.

\begin{cousinweb}
\textbf{\mr{प्रत्यक्षदर्शी}} (\emph{pratyakṣadarśī}) is a tatsama, taken from Sanskrit \emph{pratyakṣadarśin}, 'one who has seen with his own eyes' (Monier-Williams).
\end{cousinweb}

\begin{grammarlens}[title={What you've built}]
One more word for telling what happened.
\end{grammarlens}

\subsection*{Your turn}
\begin{itemize}
\raggedright
  \item \emph{Say it:} \emph{pratyakṣadarśī}
  \item \emph{Say it:} \emph{pratyakṣadarśī}, once more
  \item \emph{Say it:} \emph{ekā pratyakṣadarśīne polisānnā sagḷaṁ sāṅgitlaṁ.} — An eyewitness told the police everything.
  \item \emph{Recall:} say the Marathi for a title, then the Marathi for a front page, then say \emph{pratyakṣadarśī} again
\end{itemize}

\begin{tcolorbox}[breakable,colback=teal!4,colframe=teal!35!black,title={Before you move on}]
What does \mr{प्रत्यक्षदर्शी} mean? (\textquotedblleft{}An eyewitness\textquotedblright{}.) Say it once more.
\end{tcolorbox}
\section[\texorpdfstring{\mr{माहितीपट} (māhitīpaṭ)}{माहितीपट (māhitīpaṭ)}]{\mr{माहितीपट} (māhitīpaṭ) — a documentary}

\label{lesson:MR-C315-mahitipat}



\begin{quote}
Before the new one: say the Marathi for a front page, then the Marathi for an eyewitness.
\end{quote}

\subsection*{You'll want to know: \mr{माहितीपट}}
\textbf{\mr{माहितीपट}} — \emph{māhitīpaṭ} — \textquotedblleft{}a documentary\textquotedblright{}.

\begin{quote}
\mr{आम्ही} \mr{वाघांवरचा} \mr{एक} \mr{माहितीपट} \mr{पाहिला}.
\end{quote}

\begin{quote}
\emph{āmhī vāghāṁvarcā ek māhitīpaṭ pāhilā.}
\end{quote}

\begin{quote}
We watched a documentary about tigers.
\end{quote}

Masculine.

\mr{पट}, a screen, is masculine, and so is the compound; compare \mr{चित्रपट}, a film.

\begin{cousinweb}
\textbf{\mr{माहितीपट}} (\emph{māhitīpaṭ}) is a modern compound built from \emph{māhitī}, 'information', and Sanskrit \emph{paṭa}, 'a cloth, a screen' (Monier-Williams), on the model of \emph{citrapaṭ}, 'a film'.
\end{cousinweb}

\begin{grammarlens}[title={What you've built}]
One more word for telling what happened.
\end{grammarlens}

\subsection*{Your turn}
\begin{itemize}
\raggedright
  \item \emph{Say it:} \emph{māhitīpaṭ}
  \item \emph{Say it:} \emph{māhitīpaṭ}, once more
  \item \emph{Say it:} \emph{āmhī vāghāṁvarcā ek māhitīpaṭ pāhilā.} — We watched a documentary about tigers.
  \item \emph{Recall:} say the Marathi for a front page, then the Marathi for an eyewitness, then say \emph{māhitīpaṭ} again
\end{itemize}

\begin{tcolorbox}[breakable,colback=teal!4,colframe=teal!35!black,title={Before you move on}]
What does \mr{माहितीपट} mean? (\textquotedblleft{}A documentary\textquotedblright{}.) Say it once more.
\end{tcolorbox}
\section[\texorpdfstring{\mr{अंक} (aṅk)}{अंक (aṅk)}]{\mr{अंक} (aṅk) — a magazine issue; an issue}

\label{lesson:MR-C315-ank}



\begin{quote}
Before the new one: say the Marathi for an eyewitness, then the Marathi for a documentary.
\end{quote}

\subsection*{You'll want to know: \mr{अंक}}
\textbf{\mr{अंक}} — \emph{aṅk} — \textquotedblleft{}a magazine issue; an issue\textquotedblright{}.

\begin{quote}
\mr{मासिकाचा} \mr{नवा} \mr{अंक} \mr{मला} \mr{खूप} \mr{आवडला}.
\end{quote}

\begin{quote}
\emph{māsikācā navā aṅk malā khūp āvaḍlā.}
\end{quote}

\begin{quote}
I really liked the magazine's new issue.
\end{quote}

Masculine.

It also means a numeral: \mr{अंक} \mr{म्हणजे} \mr{आकडा}.

\begin{cousinweb}
\textbf{\mr{अंक}} (\emph{aṅk}) is a tatsama, taken from Sanskrit \emph{aṅka}, 'a hook, a mark, a numeral' (Monier-Williams); Marathi also names a numbered issue with it.
\end{cousinweb}

\begin{grammarlens}[title={What you've built}]
One more word for telling what happened.
\end{grammarlens}

\subsection*{Your turn}
\begin{itemize}
\raggedright
  \item \emph{Say it:} \emph{aṅk}
  \item \emph{Say it:} \emph{aṅk}, once more
  \item \emph{Say it:} \emph{māsikācā navā aṅk malā khūp āvaḍlā.} — I really liked the magazine's new issue.
  \item \emph{Recall:} say the Marathi for an eyewitness, then the Marathi for a documentary, then say \emph{aṅk} again
\end{itemize}

\begin{tcolorbox}[breakable,colback=teal!4,colframe=teal!35!black,title={Before you move on}]
What does \mr{अंक} mean? (\textquotedblleft{}A magazine issue; an issue\textquotedblright{}.) Say it once more.
\end{tcolorbox}
\section[\texorpdfstring{\mr{नोंदवणे} (nondavṇe)}{नोंदवणे (nondavṇe)}]{\mr{नोंदवणे} (nondavṇe) — to register; to put on record}

\label{lesson:MR-C315-nondavne}



\begin{quote}
Before the new one: say the Marathi for a documentary, then the Marathi for a magazine issue.
\end{quote}

\subsection*{You'll want to know: \mr{नोंदवणे}}
\textbf{\mr{नोंदवणे}} — \emph{nondavṇe} — \textquotedblleft{}to register; to put on record\textquotedblright{}.

\begin{quote}
\mr{हवालदाराने} \mr{माझी} \mr{तक्रार} \mr{नोंदवली}.
\end{quote}

\begin{quote}
\emph{havāldārāne mājhī takrār nondavlī.}
\end{quote}

\begin{quote}
The constable registered my complaint.
\end{quote}

Past, with -\mr{ने} (he …): \mr{त्याने} \mr{नोंदवलं}, \emph{tyāne nondavlaṁ}.

\mr{तक्रार} \mr{नोंदवणे}, to lodge a complaint; \mr{मत} \mr{नोंदवणे}, to put an opinion on record.

\begin{cousinweb}
\textbf{\mr{नोंदवणे}} (\emph{nondavṇe}) is inherited: CDIAL lists Marathi \emph{nõdṇẽ}, 'to copy and enter a document', under \emph{*nibandhati}, 'binds on; writes down' (CDIAL 7223).
\end{cousinweb}

\begin{grammarlens}[title={What you've built}]
One more word for telling what happened.
\end{grammarlens}

\subsection*{Your turn}
\begin{itemize}
\raggedright
  \item \emph{Say it:} \emph{nondavṇe}
  \item \emph{Say it:} \emph{nondavṇe}, once more
  \item \emph{Say it:} \emph{havāldārāne mājhī takrār nondavlī.} — The constable registered my complaint.
  \item \emph{Recall:} say the Marathi for a documentary, then the Marathi for a magazine issue, then say \emph{nondavṇe} again
\end{itemize}

\begin{tcolorbox}[breakable,colback=teal!4,colframe=teal!35!black,title={Before you move on}]
What does \mr{नोंदवणे} mean? (\textquotedblleft{}To register; to put on record\textquotedblright{}.) Say it once more.
\end{tcolorbox}
\section[\texorpdfstring{\mr{भटकणे} (bhaṭakṇe)}{भटकणे (bhaṭakṇe)}]{\mr{भटकणे} (bhaṭakṇe) — to wander}

\label{lesson:MR-C315-bhatakne}



\begin{quote}
Before the new one: say the Marathi for a magazine issue, then the Marathi for to register.
\end{quote}

\subsection*{You'll want to know: \mr{भटकणे}}
\textbf{\mr{भटकणे}} — \emph{bhaṭakṇe} — \textquotedblleft{}to wander\textquotedblright{}.

\begin{quote}
\mr{रस्ता} \mr{चुकल्यावर} \mr{आम्ही} \mr{दोन} \mr{तास} \mr{जंगलात} \mr{भटकलो}.
\end{quote}

\begin{quote}
\emph{rastā cukalyāvar āmhī don tās jaṅglāt bhaṭaklo.}
\end{quote}

\begin{quote}
After losing the way we wandered in the forest for two hours.
\end{quote}

Past (he …): \mr{तो} \mr{भटकला}, \emph{to bhaṭaklā}.

Intransitive; \mr{भटकी} \mr{कुत्री}, stray dogs.

\begin{cousinweb}
\textbf{\mr{भटकणे}} (\emph{bhaṭakṇe}) is inherited: CDIAL lists Marathi \emph{bhaṭakṇe}, 'to wander', under \emph{*bhaṭ}, 'a sudden movement' (CDIAL 9365).
\end{cousinweb}

\begin{grammarlens}[title={What you've built}]
One more word for telling what happened.
\end{grammarlens}

\subsection*{Your turn}
\begin{itemize}
\raggedright
  \item \emph{Say it:} \emph{bhaṭakṇe}
  \item \emph{Say it:} \emph{bhaṭakṇe}, once more
  \item \emph{Say it:} \emph{rastā cukalyāvar āmhī don tās jaṅglāt bhaṭaklo.} — After losing the way we wandered in the forest for two hours.
  \item \emph{Recall:} say the Marathi for a magazine issue, then the Marathi for to register, then say \emph{bhaṭakṇe} again
\end{itemize}

\begin{tcolorbox}[breakable,colback=teal!4,colframe=teal!35!black,title={Before you move on}]
What does \mr{भटकणे} mean? (\textquotedblleft{}To wander\textquotedblright{}.) Say it once more.
\end{tcolorbox}
